\conclusion

Ce projet de fin d'études avait pour ambition de concevoir et développer une application mobile innovante de gestion d'événements interactifs, répondant aux défis contemporains de l'éducation numérique. À travers cette démarche, nous avons réussi à créer une solution technologique complète qui transforme l'organisation et la gestion des événements éducatifs.

\textbf{Objectifs atteints :} L'application développée répond pleinement aux objectifs fixés avec 16 interfaces spécialisées organisées en cinq catégories fonctionnelles (Authentification, Gestion d'événements, Modules interactifs, Participation, et Gestion des participants). Cette architecture offre une expérience utilisateur intuitive permettant aux organisateurs de gérer efficacement des événements interactifs, tandis que les participants bénéficient d'outils d'engagement diversifiés (sondages dynamiques, quiz gamifiés, chat temps réel, sessions Q\&A).

\textbf{Innovations apportées :} L'innovation majeure réside dans l'adoption de Supabase comme Backend-as-a-Service (BaaS), permettant de développer rapidement des fonctionnalités temps réel robustes avec une base de données PostgreSQL performante et un système d'authentification sécurisé. Les souscriptions temps réel de Supabase révolutionnent l'expérience utilisateur en permettant la synchronisation instantanée des données entre participants, créant un environnement d'apprentissage véritablement interactif.

Ce travail démontre qu'une approche technologique réfléchie, combinée à une compréhension des enjeux pédagogiques, peut produire des solutions innovantes capables de transformer l'expérience d'apprentissage. L'application développée constitue un pas significatif vers une éducation plus interactive et efficace, ouvrant des perspectives d'évolution prometteuses incluant l'intelligence artificielle, les analytics avancés, et l'extension multi-plateforme.